\title{Group Sequence Policy Optimization}

\begin{abstract}
We present Group Sequence Policy Optimization (GSPO), a robust, efficient, and high-performing reinforcement learning approach designed for large language model training. Departing from earlier methods that rely on token-level importance ratios, GSPO calculates the importance ratio using sequence-level likelihood, enabling clipping, rewarding, and optimization at the sequence granularity. Our findings show that GSPO not only delivers better training efficiency and results than the GRPO method but also enhances the stability of Mixture-of-Experts (MoE) RL training and offers possibilities for streamlining RL system architectures. These advantages of GSPO have played a significant role in the advances seen in the recent Qwen3 models.
\end{abstract}

\section{Introduction}
\label{sec:intro}

Reinforcement learning (RL) has become a key framework for advancing the capabilities of language models \citep{o1,dpsk-r1,qwq32b,qwen3}. Leveraging RL at scale enables these models to address complex challenges—like high-level mathematics and programming—by engaging in more extensive and nuanced reasoning.

Ensuring consistent and resilient training dynamics is essential for effectively scaling RL with increased computational resources. Nevertheless, leading RL methods such as GRPO \citep{grpo} face pronounced stability challenges during the training of extremely large language models, frequently leading to dramatic and unrecoverable model failures \citep{qwen3, minimax-m1}. Such instability presents a significant obstacle to advancing the performance limits of language models through extended RL-based training.

In this study, we demonstrate that the core issue underlying GRPO’s instability arises from the incorrect utilization and failure of importance sampling weights within its algorithmic structure. This flaw leads to significant high-variance noise during training, which grows as response length increases and is exacerbated by the use of clipping, eventually causing the model to collapse.

In response to these fundamental challenges, we introduce \textbf{Group Sequence Policy Optimization (GSPO)}, a novel reinforcement learning algorithm designed for training large language models. GSPO’s central contribution is a principled formulation of the importance ratio using sequence likelihood, in line with established importance sampling principles \citep{click}. Moreover, GSPO leverages normalized rewards—interpreted as advantages—across multiple generated responses per query, thereby promoting coherence between sequence-level rewards and the optimization objective.

Our experimental results clearly indicate that GSPO outperforms GRPO in terms of training stability, efficiency, and overall effectiveness. Notably, GSPO intrinsically addresses the stability issues associated with reinforcement learning for large Mixture-of-Experts (MoE) models, thus obviating the reliance on intricate stabilization techniques and offering the prospect of streamlining RL system design. These advantages of GSPO were instrumental in achieving the marked performance gains observed in the recent Qwen3 models. We anticipate that GSPO will serve as a dependable and scalable methodological framework, supporting ongoing progress in large-scale RL training for language models.

\section{Preliminaries}

\paragraph{Notation}
Throughout this work, we refer to an autoregressive language model with parameters $\theta$ as a policy, denoted by $\pi_\theta$. The variable $x$ represents a query, while the entire set of queries is indicated by $\mathcal{D}$. For any response $y$ generated in reply to a query $x$, the probability assigned by the policy $\pi_\theta$ is written as $\pi_\theta (y | x)=\prod_{t=1}^{|y|} \pi_\theta (y_t | x, y_{<t} )$, where $|y|$ specifies the token count in $y$. Additionally, the tuple $(x, y)$—comprising a query and its corresponding response—can be evaluated using a verifier function $r$, which produces a reward value $r(x, y) \in [0, 1]$.

\paragraph{Proximal Policy Optimization (PPO)} Leveraging trajectories sampled under the previous policy $\pi_{\theta_\text{old}}$, PPO \citep{ppo} restricts updates to remain close to the preceding policy by incorporating a clipping strategy. In particular, PPO optimizes the following surrogate objective during policy improvement (the KL regularization term is excluded in this discussion, as it falls outside the main scope of this work):

\begin{align}
\mathcal{J}_\text{PPO}(\theta) = \mathbb{E}_{ x \sim \mathcal{D},\, y \sim \pi_{\theta_\text{old}}( \cdot | x) }
\left[ \frac{1}{|y|} \sum_{t=1}^{|y|} 
\min \left( w_{t}(\theta) \widehat{A}_{t},  \, \mathrm{clip} \left( w_{t}(\theta), 1 - {\varepsilon}, 1 + {\varepsilon}\right) \widehat{A}_{t} \right)
\right],
\end{align}

Here, the importance ratio for token $y_t$ is given by $
w_{t}(\theta) = \frac{ \pi_{\theta} (y_{t} | x, y_{<t}) }{ \pi_{\theta_\text{old}} (y_{t} | x,y_{<t})} $. The advantage $\widehat{A}_{t}$ corresponding to $y_t$ is approximated through a separate value model, while $\varepsilon$ denotes the threshold used for clipping the importance ratios.

A fundamental practical limitation of PPO is its strong dependence on the value model. In particular, the value model is typically comparable in scale to the policy model, which leads to significant demands on memory and computation. Additionally, the success of the algorithm is contingent on having accurate value estimates. Not only is it intrinsically difficult to obtain a dependable value model, but making this model effective for generating longer outputs and solving more sophisticated problems is an even more formidable obstacle.

\paragraph{Group Relative Policy Optimization (GRPO)} GRPO \citep{grpo} eliminates the reliance on a value model by evaluating the comparative advantage of each answer among a set of responses provided for the same query. In particular, GRPO seeks to maximize the following objective:

\begin{align}
\label{equ:grpo}
\mathcal{J}_\text{GRPO}(\theta) = \mathbb{E}_{ x \sim \mathcal{D},\, \{y_i\}_{i=1}^G \sim \pi_{\theta_\text{old}}( \cdot | x) }
\left[ \frac{1}{G} \sum_{i=1}^{G} \frac{1}{|y_i|} \sum_{t=1}^{|y_i|} 
\min \left( w_{i,t}(\theta) \widehat{A}_{i,t},  \, \mathrm{clip} \left( w_{i,t}(\theta), 1 - {\varepsilon}, 1 + {\varepsilon}\right) \widehat{A}_{i,t} \right)
\right],
\end{align}

Here, $G$ represents the total number of responses produced for each input $x$ (commonly referred to as the group size). The importance ratio $w_{i,t}(\theta)$ and the estimated advantage $\widehat{A}_{i,t}$ for token $y_{i,t}$ are defined as follows:

\begin{align}
    w_{i,t}(\theta)=\frac{ \pi_{\theta} (y_{i,t} | x, y_{i,<t}) }{ \pi_{\theta_\text{old}} (y_{i,t} | x,y_{i,<t})},\quad\ 
    \widehat{A}_{i,t} = \widehat{A}_{i} = \frac{r(x, y_i) - \mathrm{mean} \left( \{ r(x, y_i) \}_{i=1}^G \right) }{ \mathrm{std} \left( \{ r(x, y_i) \}_{i=1}^G \right) },
\end{align}

respectively, with each token within $y_i$ possessing the identical advantage value $\widehat{A}_{i}$.

\section{Motivation}
\label{sec:motivation}

As models become larger, incorporate greater sparsity (such as with Mixture-of-Experts architectures), and generate longer responses, it becomes essential to employ large rollout batch sizes to fully leverage hardware capabilities during RL training. For the purpose of enhancing sample efficiency, practitioners typically divide this extensive batch of rollout samples into several mini-batches used during gradient computation. This practice, however, leads to an off-policy learning context, as the collected responses $y$ originate from a preceding policy $\pi_{\theta_\text{old}}$, rather than from the updated policy $\pi_\theta$ under current optimization. This scenario underpins the rationale for incorporating clipping mechanisms in both PPO and GRPO, which are designed to mitigate the influence of samples that are too heavily "off-policy" during the estimation of gradients.

Although approaches such as clipping attempt to address off-policy divergence, we observe that GRPO suffers from a deeper limitation: \textit{the objective itself lacks well-posedness}. This flaw is especially pronounced during the training of large-scale models on tasks requiring extended responses, where it can result in severe model degradation. The core of the problem lies in an incorrect usage of importance sampling weights within the GRPO objective. Fundamentally, importance sampling is designed to estimate the expected value of a function $f$ with respect to a target distribution $\pi_\text{tar}$ by applying appropriate weighting to samples drawn from a behavior distribution $\pi_\text{beh}$:

\begin{align}
\mathbb{E}_{z \sim \pi_\text{tar}} \left[ f(z) \right] = \mathbb{E}_{z \sim \pi_\text{beh}} \left[ \frac{ \pi_\text{tar} (z) }{ \pi_\text{beh} (z)} \, f(z) \right].
\end{align}

Importantly, the effectiveness of correcting distributional discrepancies using the importance weight $\frac{ \pi_\text{tar} (z) }{ \pi_\text{beh} (z)}$ hinges on taking an average across a large number of samples ($N \gg 1$) drawn from the behavior distribution $\pi_\text{beh}$.

On the other hand, GRPO incorporates the importance weight $\frac{ \pi_{\theta} (y_{i,t} | x, y_{i,<t}) }{ \pi_{\theta_\text{old}} (y_{i,t} | x,y_{i,<t})}$ at each individual token position $t$. Because this importance weight relies solely on a single observed token $y_{i,t}$, sampled from the corresponding next-token distribution $\pi_{\theta_\text{old}}( \cdot | x, y_{i,<t})$, it does not effectively correct for distribution mismatch. Rather, it adds considerable variance to the gradient estimates during training. This high-variance noise accumulates across tokens in long sequences, with the issue further intensified by the application of the clipping strategy. Our empirical results indicate that such dynamics can precipitate irreversible model collapse: once collapse takes place, attempts to recover by restarting training—from earlier checkpoints, through careful adjustment of hyperparameters (such as clipping thresholds), by increasing sequence length, or by altering RL sampling strategies—consistently fail to restore model performance.

This observation highlights a key flaw within the design of GRPO. Specifically, the inadequacy of token-level importance weighting underscores an essential concept: \textbf{the granularity of the optimization objective must align with that of the reward}. Given that the reward is allocated for the complete sequence, implementing off-policy correction at the level of individual tokens proves to be unsuitable. This realization leads us to abandon the token-level objective in favor of leveraging importance weighting and conducting optimization at the \textit{sequence level} instead.

\section{Algorithm}

\subsection{GSPO: Group Sequence Policy Optimization}

Although the token-level importance weight $\frac{ \pi_{\theta} (y_{i,t} | x, y_{i,<t}) }{ \pi_{\theta_\text{old}} (y_{i,t} | x,y_{i,<t})}$ presents challenges within GRPO, our analysis indicates that the \textit{sequence-level} importance weight $\frac{ \pi_{\theta} (y | x) }{ \pi_{\theta_\text{old}} (y | x)}$ offers a well-defined theoretical interpretation in language generation settings. Specifically, it quantifies the divergence between a response $y$ generated under $\pi_{\theta_\text{old}} (\cdot | x)$ and the target distribution $\pi_{\theta} (\cdot | x)$. This measure is inherently compatible with sequence-level rewards and effectively guides the operation of the clipping mechanism.

Drawing from this basic insight, we introduce the \textbf{Group Sequence Policy Optimization (GSPO)} algorithm. The GSPO method utilizes the subsequent optimization objective defined at the sequence level:

\begin{align}
\mathcal{J}_\text{GSPO} (\theta) =
\mathbb{E}_{ x \sim \mathcal{D},\, \{y_i\}_{i=1}^G \sim \pi_{\theta_\text{old}}( \cdot | x) }
\left[ 
\frac{1}{G} \sum_{i=1}^{G}
\min \left( s_{i}(\theta)  \widehat{A}_{i},  \, \mathrm{clip} \left( s_{i}(\theta), 1 - {\varepsilon}, 1 + {\varepsilon} \right) \widehat{A}_{i} \right) 
\right],
\label{equ:gspo}
\end{align}

in which we utilize the advantage estimation based on groupings:

\begin{align}
\widehat{A}_{i} = \frac{r(x, y_i) - \mathrm{mean} \left( \{ r(x, y_i) \}_{i=1}^G \right) }{ \mathrm{std} \left( \{ r(x, y_i) \}_{i=1}^G \right) },
\end{align}

and establish the significance ratio $s_{i}(\theta)$ grounded in the probability of the sequence \citep{click}:

\begin{align}
s_{i}(\theta) = \left( \frac{ \pi_{\theta} (y_i | x) }{ \pi_{\theta_\text{old}} (y_i | x)} \right)^{\frac{1}{|y_i|}}
=
\exp \left( \frac{1}{|y_i|} \sum_{t=1}^{|y_i|} \log \frac{ \pi_{\theta} (y_{i,t} | x, y_{i,<t}) }{ \pi_{\theta_\text{old}} (y_{i,t} | x,y_{i,<t})} \right).
\end{align}

Consequently, GSPO performs clipping at the level of entire responses, rather than on a token-by-token basis, in order to remove excessively ``off-policy'' samples during gradient computation. This approach aligns the reward computation and the optimization process at the sequence level. Additionally, we implement length normalization for $s_{i}(\theta)$, aiming to minimize variance and maintain $s_{i}(\theta)$ within a consistent numerical interval. Without normalization, significant changes in a small subset of token probabilities could cause large swings in the sequence-level importance ratio, necessitating different clipping ranges for responses depending on their lengths. Furthermore, it is important to highlight that, because GSPO and earlier methods (such as GRPO) define importance ratios differently, the resulting clipping ranges can vary considerably, often by an order of magnitude.

\subsection{Gradient Analysis}
\label{subsec:gradient}

The gradient of the GSPO objective can be calculated in the following manner (for simplicity, the clipping step is excluded):

\begin{align}
\nabla_{\theta} \mathcal{J}_\text{GSPO} (\theta)
=&\ 
\nabla_{\theta} \mathbb{E}_{ x \sim \mathcal{D},\, \{y_i\}_{i=1}^G \sim \pi_{\theta_\text{old}}( \cdot | x) }
\left[ 
\frac{1}{G} \sum_{i=1}^{G} s_{i}(\theta) \widehat{A}_{i}
\right] \\
= &\ 
\mathbb{E}_{ x \sim \mathcal{D},\, \{y_i\}_{i=1}^G \sim \pi_{\theta_\text{old}}( \cdot | x) }
\left[ 
\frac{1}{G} \sum_{i=1}^{G}
s_{i}(\theta) \widehat{A}_{i}
\cdot \nabla_{\theta} \log s_{i}(\theta)
\right] \\
=&\ 
\mathbb{E}_{ x \sim \mathcal{D},\, \{y_i\}_{i=1}^G \sim \pi_{\theta_\text{old}}( \cdot | x) }
\left[ 
\frac{1}{G} \sum_{i=1}^{G} 
\left( \frac{ \pi_{\theta} (y_i | x) }{ \pi_{\theta_\text{old}} (y_i | x)} \right)^{\frac{1}{|y_i|}} \widehat{A}_{i} 
\cdot \frac{1}{|y_i|} \sum_{t=1}^{|y_i|} \nabla_{\theta} \log \pi_{\theta} (y_{i,t} | x, y_{i,<t}) 
\right].
\label{equ:gspo_grad}
\end{align}

To facilitate comparison, the gradient for the GRPO objective can be written as (recall that $\widehat{A}_{i,t} = \widehat{A}_{i}$):

\begin{align}
\nabla_{\theta} \mathcal{J}_\text{GRPO} (\theta) 
=&\ 
\nabla_{\theta} \mathbb{E}_{ x \sim \mathcal{D},\, \{y_i\}_{i=1}^G \sim \pi_{\theta_\text{old}}( \cdot | x) }
\left[ \frac{1}{G} \sum_{i=1}^{G} \frac{1}{|y_i|} \sum_{t=1}^{|y_i|} 
w_{i,t}(\theta) \widehat{A}_{i,t}
\right] \\
=&\
\mathbb{E}_{ x \sim \mathcal{D},\, \{y_i\}_{i=1}^G \sim \pi_{\theta_\text{old}}( \cdot | x) }
\left[ \frac{1}{G} \sum_{i=1}^{G} \widehat{A}_{i} 
\cdot \frac{1}{|y_i|} \sum_{t=1}^{|y_i|} 
\frac{ \pi_{\theta} (y_{i,t} | x, y_{i,<t}) }{ \pi_{\theta_\text{old}} (y_{i,t} | x,y_{i,<t})} 
\nabla_{\theta} \log \pi_{\theta} (y_{i,t} | x, y_{i,<t})  
\right].
\end{align}

Consequently, the principal difference between GSPO and GRPO pertains to \textit{the method by which gradients of token log likelihoods are weighted}. In the case of GRPO, each token is assigned a weight that corresponds to its ``importance weight'' $\frac{ \pi_{\theta} (y_{i,t} | x, y_{i,<t}) }{ \pi_{\theta_\text{old}} (y_{i,t} | x,y_{i,<t})}$. These importance weights, which fall within the intervals $(0, 1+\varepsilon]$ when $\widehat{A}_{i} >0$, or $[1-\varepsilon, +\infty)$ when $\widehat{A}_{i} < 0$, can differ significantly among tokens. Such variability in weighting is substantial and can accumulate over the course of training, introducing the potential for erratic or unintended behavior. On the other hand, GSPO assigns uniform weight to every token within a response, thereby removing the aforementioned source of instability inherent to GRPO.

\subsection{GSPO-token: A Token-level Objective Variant}

In situations such as multi-turn RL, there may be a need to modify the advantage at a more granular level than the entire sequence. Therefore, we propose a token-level adaptation of GSPO, referred to as \textbf{GSPO-token}, which facilitates the adjustment of advantages for individual tokens:

\begin{align}
\mathcal{J}_\text{GSPO-token}(\theta)
=
\mathbb{E}_{ x \sim \mathcal{D},\, \{y_i\}_{i=1}^G \sim \pi_{\theta_\text{old}}( \cdot | x) }
\left[ 
\frac{1}{G} \sum_{i=1}^{G} \frac{1}{|y_i|} \sum_{t=1}^{|y_i|} 
\min \left( s_{i,t}(\theta) \widehat{A}_{i,t},  \, \mathrm{clip} \left( s_{i,t}(\theta), 1 - {\varepsilon}, 1 + {\varepsilon} \right) \widehat{A}_{i,t} \right) 
\right],
\label{equ:gspo-token}
\end{align}

in which

\begin{align}
s_{i,t}(\theta) 
= \mathrm{sg} \left[ s_{i}(\theta) \right]  \cdot \frac{ \pi_{\theta} (y_{i,t} | x, y_{i,<t}) }{ \mathrm{sg} \left[ \pi_{\theta} (y_{i,t} | x, y_{i,<t}) \right] },
\end{align}

and $\mathrm{sg}[\cdot]$ refers to extracting the numerical value while halting any gradient flow, analogous to the \texttt{detach} function in PyTorch. The gradient expression for GSPO-token can be obtained as follows:

\begin{align}
\nabla_{\theta} \mathcal{J}_\text{GSPO-token}(\theta)
=&\ 
\nabla_{\theta} \mathbb{E}_{ x \sim \mathcal{D},\, \{y_i\}_{i=1}^G \sim \pi_{\theta_\text{old}}( \cdot | x) }
\left[ 
\frac{1}{G} \sum_{i=1}^{G}
\frac{1}{|y_i|} \sum_{t=1}^{|y_i|} 
s_{i,t}(\theta) \widehat{A}_{i,t}
\right] \\
=&\ 
\mathbb{E}_{ x \sim \mathcal{D},\, \{y_i\}_{i=1}^G \sim \pi_{\theta_\text{old}}( \cdot | x) }
\left[ 
\frac{1}{G} \sum_{i=1}^{G} s_i (\theta) \cdot \frac{1}{|y_i|} \sum_{t=1}^{|y_i|} 
\widehat{A}_{i,t} \frac{ \nabla_{\theta} \pi_{\theta} (y_{i,t} | x, y_{i,<t}) }{ \pi_{\theta} (y_{i,t} | x, y_{i,<t}) }
\right] \\
=&\ 
\mathbb{E}_{ x \sim \mathcal{D},\, \{y_i\}_{i=1}^G \sim \pi_{\theta_\text{old}}( \cdot | x) }
\left[ 
\frac{1}{G} \sum_{i=1}^{G}  \left( \frac{ \pi_{\theta} (y_i | x) }{ \pi_{\theta_\text{old}} (y_i | x)} \right)^{\frac{1}{|y_i|}}
\cdot \frac{1}{|y_i|} \sum_{t=1}^{|y_i|}
\widehat{A}_{i,t} \nabla_{\theta} \log \pi_{\theta} (y_{i,t} | x, y_{i,<t})  
\right].
\label{equ:gspo-token_grad}
\end{align}

Observe that the expression $\frac{ \pi_{\theta} (y_{i,t} | x, y_{i,<t}) }{ \mathrm{sg} \left[ \pi_{\theta} (y_{i,t} | x, y_{i,<t}) \right] }$ evaluates to 1, which implies that $s_{i,t}(\theta)$ is, in terms of actual value, equivalent to $s_{i}(\theta)$.
By examining Equation~\eqref{equ:gspo} alongside~\eqref{equ:gspo-token}, as well as Equation~\eqref{equ:gspo_grad} with~\eqref{equ:gspo-token_grad}, it becomes clear that GSPO-token and GSPO are exactly the same in their numerical optimization objective, clipping mechanism, and theoretical gradient, provided all token-level advantages within the sequence $y_i$ are uniformly set (that is, $\widehat{A}_{i,t} = \widehat{A}_{i}$ for each $t$). Nevertheless, GSPO-token provides the added benefit of being able to specify token-wise advantages, which introduces greater flexibility.

\section{Experiments and Discussion}

\subsection{Empirical Results}

We conduct experiments using a cold-start model that is fine-tuned from Qwen3-30B-A3B-Base, presenting both training reward trajectories and performance metrics across the AIME'24 (average Pass@1 from 32 samples), LiveCodeBench (202410-202502, average Pass@1 from 8 samples), and CodeForces (Elo Rating) evaluation suites. In our reinforcement learning setup, the rollout data in each batch is divided into four mini-batches to facilitate gradient updates. For GSPO, the clipping intervals in Equation~\eqref{equ:gspo} are assigned to 3e-4 (left) and 4e-4 (right). For comparative purposes, we benchmark against GRPO, using this method as our baseline with clipping ranges in Equation~\eqref{equ:grpo} set to 0.2 (left) and 0.27 (right); these parameters are the result of meticulous tuning for an equitable assessment. It is important to highlight that GRPO relies on the Routing Replay training approach to attain reliable MoE RL convergence—a subject we revisit in \S~\ref{subsec:routing_replay}—whereas \textit{GSPO removes the dependency on this training strategy}.

Figure~\ref{fig:results} illustrates that training with GSPO maintains stable progress at all stages. Our findings indicate that \textbf{GSPO achieves ongoing gains in performance when more training compute is allocated, the query set is periodically refreshed, and generation length is increased}. In addition, GSPO outperforms GRPO in terms of training efficiency, reaching higher levels of training accuracy and benchmark results while using identical amounts of training compute and query resources. Lastly, we have effectively leveraged GSPO in the RL training of the state-of-the-art Qwen3 models, providing strong evidence for GSPO's capability in harnessing RL scaling benefits for large language models.

\subsection{Curious Observation on Clipping Fractions}

One major difference between GSPO and GRPO lies in GSPO’s approach of clipping whole responses, as opposed to GRPO’s token-wise clipping strategy. Notably, as illustrated in Figure~\ref{fig:clipping}, there is a disparity of approximately two orders of magnitude between the proportions of clipped tokens under GSPO and GRPO. Adjusting the clipping intervals does not eliminate this significant gap. Surprisingly, although GSPO clips far more tokens—and thus utilizes a smaller portion of data for both training and gradient calculation—it nonetheless surpasses GRPO in training efficiency. This unexpected outcome, wherein a much greater proportion of clipped tokens correlates with improved training performance, suggests that the token-level gradient approximations employed by GRPO are subject to high variance and are suboptimal for effective sample utilization. Conversely, GSPO’s reliance on sequence-level gradients yields a learning signal that is both more stable and productive.

\subsection{Benefit of GSPO for MoE Training}
\label{subsec:routing_replay}

\paragraph{Background}
While reinforcement learning (RL) training of dense models generally proceeds smoothly, the inherently sparse activation mechanism in MoE models presents distinct stability concerns. Specifically, our investigation revealed that under the GRPO algorithm, MoE models are susceptible to \textit{expert-activation volatility}, which can impede the successful convergence of RL optimization. Notably, following one or several gradient update steps, the subset of experts triggered for generating an identical response often exhibits substantial variation. As an illustration, using the 48-layer Qwen3-30B-A3B-Base model, approximately 10\% of the experts engaged by the updated policy $\pi_{\theta}$—on the same rollout sample—differ from those engaged by the previous policy $\pi_{\theta_\text{old}}$ after every RL gradient step. This volatility is further exacerbated in MoE models of greater depth, resulting in highly unstable token-level importance ratios $w_{i,t}(\theta) = \frac{ \pi_{\theta} (y_{i,t} | x, y_{i,<t}) }{ \pi_{\theta_\text{old}} (y_{i,t} | x,y_{i,<t})}$, as previously addressed in \S~\ref{sec:motivation} and~\ref{subsec:gradient}. Such instability undermines the effectiveness of these ratios and disrupts the normal convergence behavior of RL training.

\paragraph{Our Previous Approach} In addressing this issue, our earlier work introduced the \textbf{Routing Replay} training strategy. The approach involves storing the experts activated by $\pi_{\theta_\text{old}}$ and subsequently ``replaying'' these identical routing configurations in $\pi_{\theta}$ during the calculation of importance ratios $w_{i,t}(\theta) = \frac{ \pi_{\theta} (y_{i,t} | x, y_{i,<t}) }{ \pi_{\theta_\text{old}} (y_{i,t} | x,y_{i,<t})}$. By ensuring that both $\pi_{\theta} (y_{i,t} | x, y_{i,<t})$ and $\pi_{\theta_\text{old}} (y_{i,t} | x,y_{i,<t})$ utilize the same network paths for each token $y_{i,t}$, this method restores stability to the importance ratios at the token level. Additionally, it facilitates consistent optimization of the activated network components throughout gradient updates. As illustrated in Figure~\ref{fig:routing_replay}, Routing Replay proves to be a vital component for achieving standard convergence behavior during GRPO training in MoE architectures.

\paragraph{Benefit of GSPO} While Routing Replay helps ensure that GRPO-based training of MoE models reaches proper convergence, its strategy of recycling routing patterns introduces extra memory demands and communication costs, potentially constraining the model's effective capacity. In comparison, as depicted in Figure~\ref{fig:results}, GSPO removes reliance on Routing Replay altogether and supports standard computation of the importance ratios $s_i(\theta)$, achieving normal convergence and stable optimization. The central advantage is that GSPO considers only the sequence-level likelihood (i.e., $\pi_{\theta} (y_i | x)$), rather than being influenced by the per-token likelihood (i.e., $\pi_{\theta} (y_{i,t} | x, y_{i,<t})$). Given that the MoE model consistently retains robust language modeling abilities, the overall sequence likelihood tends to remain steady. In essence, GSPO directly addresses the challenge of volatility in expert activation for MoE architectures, removing the necessity for intricate solutions like Routing Replay. This not only streamlines and stabilizes training, but also empowers the model to utilize its entire representational capacity without imposed limitations.

\subsection{Benefit of GSPO for RL Infrastructure}

Due to differences in numerical precision between training platforms such as Megatron and inference platforms like SGLang or vLLM, it is standard practice to rely on the training engine to recalculate the likelihoods of generated responses based on the old policy $\pi_{\theta_\text{old}}$. In contrast, GSPO leverages sequence-level likelihoods during optimization, as opposed to token-level ones, and sequence-level calculations are generally far less sensitive to minor precision mismatches. Consequently, GSPO enables the direct use of likelihoods provided by the inference engine within the optimization loop, removing the necessity for recomputation with the original training infrastructure. This approach proves especially advantageous in settings such as partial rollout, multi-turn RL, and systems where training and inference are conducted separately.

\section{Conclusion}

We introduce Group Sequence Policy Optimization (GSPO), an innovative reinforcement learning algorithm designed for large language model training. Building on the core concept of importance sampling, GSPO formulates importance ratios using sequence likelihoods, and applies sequence-level clipping, reward assignment, and optimization processes. In empirical studies, GSPO achieves greater stability, training efficiency, and overall performance when contrasted with GRPO, and is especially advantageous for large-scale reinforcement learning applications involving MoE models. This advancement underpins the marked progress realized in the newest Qwen3 models. By establishing GSPO as a scalable, foundational technique, we aim to further expand RL scaling and anticipate significant future progress in artificial intelligence capabilities.

\end{document}